% Part: sets-functions-relations
% Chapter: arithmetization
% Section: From N to Z

\documentclass[../../../include/open-logic-section]{subfiles}


\begin{document}

\olfileid{sfr}{arith}{int}
\olsection{$\Nat$ నుంచి $\Int$కు}

ఇవిగో రెండు ప్రాథమిక అవగాహనలు:
	\begin{enumerate}
		\item ప్రతి పూర్ణ సంఖ్యను $n - m$ రూపంలో రాయవచ్చు; ఇక్కడ $n, m \in\Nat$.
		\item $n - m$ అనే వ్యక్తీకరణలో సంకేతీకరించిన సమాచారాన్ని
		క్రమయుగ్మం $\tuple{n, m}$తో కూడా అంతే విధంగా సంకేతీకరించవచ్చు.
	\end{enumerate}
సహజ సంఖ్యల క్రమయుగ్మాలు $\Nat^2$ యొక్క \tetoken{మూలకాలు}{!!{element}s} అని మనకు ఇప్పటికే
తెలుసు. అంతేకాక, $\Nat$ను మనం అర్థం చేసుకున్నామని అనుకుంటున్నాం.
కాబట్టి పై రెండు అవగాహనల ఆధారంగా ఒక సరళమైన సూచన ఇదిగో:
\emph{పూర్ణ సంఖ్యలను సహజ సంఖ్యల క్రమయుగ్మాలుగా పరిగణిద్దాం}.

నిజానికి ఈ సూచన మరీ సరళమైనది. $0- 2 = 4 - 6$ కావాలని మనం స్పష్టంగా
కోరుకుంటాం. కానీ $\tuple{0, 2 }\neq \tuple{4, 6}$ అనేది స్పష్టం.
అందువల్ల $\Nat^2$నే పూర్ణ సంఖ్యల సమితి అని నేరుగా చెప్పలేం.

ముందున్న సమస్యను సాధారణీకరిస్తే, మనకు కావలసింది ఇది:
	$$a - b = c - d \text{ iff }a + d = c + b$$
(పూర్ణ సంఖ్యలు ఇలాగే \emph{ప్రవర్తించాలి} అనేది స్పష్టమే: రెండు
వైపులకూ $b$, $d$లను కలపండి.) ఈ ప్రవర్తనకు హామీ ఇచ్చే సులభమైన మార్గం,
క్రమయుగ్మాల మధ్య $\Intequiv$ అనే తుల్యతా సంబంధాన్ని ఇలా నిర్వచించడం:
	$$\tuple{ a, b } \Intequiv \tuple{c, d}\text{ iff }a + d = c + b$$
ఇప్పుడు ఇది తుల్యతా సంబంధమని చూపాలి.
\begin{prop} $\Intequiv$ ఒక తుల్యతా సంబంధం.
\end{prop}
	\begin{proof}
	$\Intequiv$ స్వావర్తన, సౌష్ఠవ, సంక్రామక ధర్మాలు కలిగి ఉందని
	చూపాలి.

	\emph{స్వావర్తన ధర్మం:} $a + b = b + a$ కనుక
$\tuple{a, b} \Intequiv \tuple{a, b}$ అనేది స్పష్టం.

	\emph{సౌష్ఠవ ధర్మం:} $\tuple{a, b} \Intequiv \tuple{c, d}$
అనుకుందాం; అందువల్ల $a + d = c + b$. అప్పుడు $c + b = a + d$;
కాబట్టి $\tuple{c, d} \Intequiv \tuple{a, b}$.

	\emph{సంక్రామక ధర్మం:} $\tuple{a, b} \Intequiv \tuple{c,
d}\Intequiv \tuple{m, n}$ అనుకుందాం. అప్పుడు $a + d = c + b$ మరియు
$c + n = m + d$. కనుక $a + d + c + n = c + b + m + d$; అందువల్ల
$a + n = m + b$. కాబట్టి $\tuple{a, b } \Intequiv \tuple{m, n}$.
\end{proof}

ఇప్పుడు ఈ తుల్యతా సంబంధాన్ని ఉపయోగించి తుల్యతా వర్గాలను తీసుకోవచ్చు.

\begin{defn}
		పూర్ణ సంఖ్యలు అంటే $\Intequiv$ కింద సహజ సంఖ్యల క్రమయుగ్మాల
		తుల్యతా వర్గాలు; అంటే,
$\Int = \equivclass{\Nat^2}{\Intequiv}$.
\end{defn}

ఈ నిర్దేశిత నిర్వచనం గురించి అనేక భిన్నమైన \emph{తాత్త్విక}
ప్రతిస్పందనలు ఉండవచ్చు. అయితే వాటిని పరిశీలించే ముందు కొన్ని సాంకేతిక
వివరాలను కొనసాగించడం ఉపయోగకరం.

పూర్ణ సంఖ్యలు ఏమిటో చెప్పిన తరువాత, వాటిపై ప్రాథమిక ప్రమేయాలు,
సంబంధాలు నిర్వచించాలి. $\tuple{m, n}$ను \tetoken{ఒక మూలకంగా}{!!a{element}} కలిగిన,
$\Intequiv$ కింద తుల్యతా వర్గాన్ని $\equivrep{m, n}{\Intequiv}$గా
రాద్దాం.\footnote{గమనిక: \olref[sfr][rel][eqv]{def:equivalenceclass}లో
ప్రవేశపెట్టిన సంజ్ఞను వాడితే, ఇదే దానిని
$\equivrep{\tuple{m,n}}{\Intequiv}$గా రాసేవాళ్లం. కానీ అది చదవడానికి
కొంచెం కష్టం.} అంటే:
	$$\equivrep{m, n}{\Intequiv} = \Setabs{\tuple{a, b} \in \Nat^2}{\tuple{a, b}\Intequiv \tuple{m, n}}$$
ఇప్పుడు ఈ నిర్వచనాలను ఇస్తాం:
	\begin{align*}
		\equivrep{a, b}{\Intequiv} + \equivrep{c, d}{\Intequiv} &= \equivrep{a + c, b + d}{\Intequiv}\\
		\equivrep{a, b}{\Intequiv} \times \equivrep{c, d}{\Intequiv} &= \equivrep{a c + b  d, a  d + b c}{\Intequiv}\\
		\equivrep{a, b}{\Intequiv} \leq \equivrep{c, d}{\Intequiv} &\text{ iff }a + d \leq b + c
	\end{align*}
(సాధారణంగా చేసే విధంగానే, సూత్రాలు చదవడానికి సులభంగా ఉండేందుకు
`$ab$'ను `$(a \times b)$'కు సంక్షిప్తంగా వాడుతున్నాను.) ఇప్పుడు ఈ
నిర్వచనాలు అవి \emph{ప్రవర్తించాల్సిన} విధంగానే ప్రవర్తిస్తాయని
నిర్ధారించాలి. దాని అర్థాన్ని పూర్తిగా చెప్పి తనిఖీ చేయడం శ్రమతో కూడిన
పని; వివరాలను \olref[check]{sec}కు వదిలేస్తాం. సంక్షిప్తంగా చెప్పాలంటే:
అన్నీ పనిచేస్తాయి!

చివరగా ఇంకొక విషయం మిగిలింది. సహజ సంఖ్యలను ఉపయోగించి పూర్ణ సంఖ్యలను
నిర్మించాం. అయితే దాని వల్ల సహజ సంఖ్యలు \emph{స్వయంగా పూర్ణ సంఖ్యలు
కావు}. దీని తాత్త్విక ప్రాముఖ్యతను \olref[ref]{sec}లో మళ్లీ
పరిశీలిస్తాం. పూర్తిగా సాంకేతిక దృష్టిలో మాత్రం, సహజ సంఖ్యలను పూర్ణ
సంఖ్యలుగా పరిగణించడానికి ఒక మార్గం కావాలి. ఆలోచన సులభం: ప్రతి
$n \in \Nat$కు $n_\Int = \equivrep{n, 0}{\Intequiv}$ అని
నిర్దేశిద్దాం. ఈ నిర్వచనం సరిగ్గా ప్రవర్తిస్తుందని, అంటే ఏ
$m, n \in \Nat$కైనా కింది సమీకరణాలు వర్తిస్తాయని నిర్ధారించాలి:
\begin{align*}
	(m + n)_\Int &= m_\Int + n_\Int\\
	(m \times n)_\Int &= m_\Int \times n_\Int\\
	m \leq n &\liff m_\Int \leq n_\Int
\end{align*}
ఇదంతా సూటిగానే ఉంటుంది. ఉదాహరణకు, వీటిలో రెండవది వర్తిస్తుందని
చూపడానికి సహజ సంఖ్యల ప్రవర్తనను ఉపయోగించి ఇలా తర్కించవచ్చు:
%\begin{proof}
%	సహజ సంఖ్యల ప్రవర్తనను ఉపయోగిస్తే, స్పష్టంగా:
	\begin{align*}
%		(m + n)_\Int  = \equivrep{m + n, 0}{\Intequiv} = \equivrep{m + n, 0 + 0}{\Intequiv} = \equivrep{m, 0}{\Intequiv} + \equivrep{n, 0}{\Intequiv} = m_\Int + n_\Int\\
		(m \times n)_\Int  &= \equivrep{m \times n, 0}{\Intequiv} \\
		&= \equivrep{m \times n + 0 \times 0, m \times 0 + 0 \times n}{\Intequiv} \\
		&= \equivrep{m, 0}{\Intequiv} \times \equivrep{n, 0}{\Intequiv} \\
		&= m_\Int \times n_\Int
	\end{align*}
%\end{proof}
మిగిలిన రెండు షరతులు వర్తిస్తాయని నిర్ధారించడాన్ని సాధనగా
వదిలేస్తాం.
\begin{prob}
	ఏ $m, n \in \Nat$కైనా $(m + n)_\Int = m_\Int + n_\Int$ మరియు
$m \leq n \liff m_\Int \leq n_\Int$ అని చూపండి.
\end{prob}
\end{document}
